\PassOptionsToPackage{unicode}{hyperref}
\PassOptionsToPackage{hyphens}{url}
\PassOptionsToPackage{dvipsnames,svgnames*,x11names*}{xcolor}
\documentclass[
  11pt,
]{article}
\usepackage{amsmath,amssymb}
\usepackage{lmodern}
\usepackage{iftex}
\ifPDFTeX
  \usepackage[T1]{fontenc}
  \usepackage[utf8]{inputenc}
  \usepackage{textcomp}
\else
  \usepackage{unicode-math}
  \defaultfontfeatures{Scale=MatchLowercase}
  \defaultfontfeatures[\rmfamily]{Ligatures=TeX,Scale=1}
\fi
\IfFileExists{upquote.sty}{\usepackage{upquote}}{}
\IfFileExists{microtype.sty}{
  \usepackage[]{microtype}
  \UseMicrotypeSet[protrusion]{basicmath}
}{}
\makeatletter
\@ifundefined{KOMAClassName}{
  \IfFileExists{parskip.sty}{
    \usepackage{parskip}
  }{
    \setlength{\parindent}{0pt}
    \setlength{\parskip}{6pt plus 2pt minus 1pt}}
}{
  \KOMAoptions{parskip=half}}
\makeatother
\usepackage{xcolor}
\IfFileExists{xurl.sty}{\usepackage{xurl}}{}
\IfFileExists{bookmark.sty}{\usepackage{bookmark}}{\usepackage{hyperref}}
\hypersetup{
  colorlinks=true,
  linkcolor={blue},
  filecolor={Maroon},
  citecolor={Blue},
  urlcolor={blue},
  pdfcreator={LaTeX via pandoc}}
\urlstyle{same}
\usepackage[margin=2cm]{geometry}
\usepackage{color}
\usepackage{fancyvrb}
\newcommand{\VerbBar}{|}
\newcommand{\VERB}{\Verb[commandchars=\\\{\}]}
\DefineVerbatimEnvironment{Highlighting}{Verbatim}{commandchars=\\\{\}}
\newenvironment{Shaded}{}{}
\newcommand{\AlertTok}[1]{\textcolor[rgb]{1.00,0.00,0.00}{\textbf{#1}}}
\newcommand{\AnnotationTok}[1]{\textcolor[rgb]{0.38,0.63,0.69}{\textbf{\textit{#1}}}}
\newcommand{\AttributeTok}[1]{\textcolor[rgb]{0.49,0.56,0.16}{#1}}
\newcommand{\BaseNTok}[1]{\textcolor[rgb]{0.25,0.63,0.44}{#1}}
\newcommand{\BuiltInTok}[1]{#1}
\newcommand{\CharTok}[1]{\textcolor[rgb]{0.25,0.44,0.63}{#1}}
\newcommand{\CommentTok}[1]{\textcolor[rgb]{0.38,0.63,0.69}{\textit{#1}}}
\newcommand{\CommentVarTok}[1]{\textcolor[rgb]{0.38,0.63,0.69}{\textbf{\textit{#1}}}}
\newcommand{\ConstantTok}[1]{\textcolor[rgb]{0.53,0.00,0.00}{#1}}
\newcommand{\ControlFlowTok}[1]{\textcolor[rgb]{0.00,0.44,0.13}{\textbf{#1}}}
\newcommand{\DataTypeTok}[1]{\textcolor[rgb]{0.56,0.13,0.00}{#1}}
\newcommand{\DecValTok}[1]{\textcolor[rgb]{0.25,0.63,0.44}{#1}}
\newcommand{\DocumentationTok}[1]{\textcolor[rgb]{0.73,0.13,0.13}{\textit{#1}}}
\newcommand{\ErrorTok}[1]{\textcolor[rgb]{1.00,0.00,0.00}{\textbf{#1}}}
\newcommand{\ExtensionTok}[1]{#1}
\newcommand{\FloatTok}[1]{\textcolor[rgb]{0.25,0.63,0.44}{#1}}
\newcommand{\FunctionTok}[1]{\textcolor[rgb]{0.02,0.16,0.49}{#1}}
\newcommand{\ImportTok}[1]{#1}
\newcommand{\InformationTok}[1]{\textcolor[rgb]{0.38,0.63,0.69}{\textbf{\textit{#1}}}}
\newcommand{\KeywordTok}[1]{\textcolor[rgb]{0.00,0.44,0.13}{\textbf{#1}}}
\newcommand{\NormalTok}[1]{#1}
\newcommand{\OperatorTok}[1]{\textcolor[rgb]{0.40,0.40,0.40}{#1}}
\newcommand{\OtherTok}[1]{\textcolor[rgb]{0.00,0.44,0.13}{#1}}
\newcommand{\PreprocessorTok}[1]{\textcolor[rgb]{0.74,0.48,0.00}{#1}}
\newcommand{\RegionMarkerTok}[1]{#1}
\newcommand{\SpecialCharTok}[1]{\textcolor[rgb]{0.25,0.44,0.63}{#1}}
\newcommand{\SpecialStringTok}[1]{\textcolor[rgb]{0.73,0.40,0.53}{#1}}
\newcommand{\StringTok}[1]{\textcolor[rgb]{0.25,0.44,0.63}{#1}}
\newcommand{\VariableTok}[1]{\textcolor[rgb]{0.10,0.09,0.49}{#1}}
\newcommand{\VerbatimStringTok}[1]{\textcolor[rgb]{0.25,0.44,0.63}{#1}}
\newcommand{\WarningTok}[1]{\textcolor[rgb]{0.38,0.63,0.69}{\textbf{\textit{#1}}}}
\usepackage{longtable,booktabs,array}
\usepackage{calc}
\usepackage{etoolbox}
\makeatletter
\patchcmd\longtable{\par}{\if@noskipsec\mbox{}\fi\par}{}{}
\makeatother
\IfFileExists{footnotehyper.sty}{\usepackage{footnotehyper}}{\usepackage{footnote}}
\makesavenoteenv{longtable}
\setlength{\emergencystretch}{3em}
\providecommand{\tightlist}{
  \setlength{\itemsep}{0pt}\setlength{\parskip}{0pt}}
\setcounter{secnumdepth}{5}
\usepackage{fvextra}
\DefineVerbatimEnvironment{Highlighting}{Verbatim}{breaklines,breakanywhere,fontsize=\footnotesize,commandchars=\\\{\}}
\RecustomVerbatimEnvironment{verbatim}{Verbatim}{breaklines,breakanywhere,fontsize=\footnotesize}
\usepackage{microtype}
\usepackage{booktabs}
\usepackage{longtable}
\usepackage{array}
\usepackage{enumitem}
\setlist[itemize]{itemsep=2pt,parsep=2pt}
\setlist[enumerate]{itemsep=2pt,parsep=2pt}
\usepackage{etoolbox}
\AtBeginEnvironment{longtable}{\footnotesize}
\AtBeginEnvironment{tabular}{\footnotesize}
\renewcommand{\arraystretch}{1.2}
\setlength{\tabcolsep}{4pt}
\usepackage{fancyhdr}
\usepackage{titlesec}
\titleformat{\section}{\Large\bfseries}{\thesection}{0.6em}{}
\titleformat{\subsection}{\large\bfseries}{\thesubsection}{0.6em}{}
\titlespacing*{\section}{0pt}{1.4em}{0.6em}
\titlespacing*{\subsection}{0pt}{1.0em}{0.4em}
\setlength{\parindent}{0pt}
\setlength{\parskip}{0.5em plus 0.1em minus 0.1em}
\pagestyle{fancy}
\fancyhf{}
\fancyhead[L]{\nouppercase{\leftmark}}
\fancyhead[R]{\thepage}
\fancyfoot[C]{overlap-calculator case studies}
\renewcommand{\headrulewidth}{0.4pt}
\renewcommand{\footrulewidth}{0pt}
\setlength{\headheight}{14pt}
\sloppy
\ifLuaTeX
  \usepackage{selnolig}
\fi

\title{Case Study 14 --- Calibration Block}
\usepackage{etoolbox}
\makeatletter
\providecommand{\subtitle}[1]{
  \apptocmd{\@title}{\par {\large #1 \par}}{}{}
}
\makeatother
\subtitle{Linear energy map, oscillator scaling, and per-band width overrides}
\author{}
\date{}

\begin{document}
\maketitle

{
\hypersetup{linkcolor=}
\setcounter{tocdepth}{2}
\tableofcontents
}
\hypertarget{case-study-14--calibration-block}{
\section{Case Study 14 --- Calibration
Block}\label{case-study-14--calibration-block}}

Apply a TOML calibration before broadening. The calibration block can
shift TD-DFT excitation energies with a linear map
\(E_{\text{cal}} = aE + b\), scale oscillator strengths
\(f_{\text{cal}} = \alpha f\), and override the broadening width or
reorganization energy per excited state.

The calibration is applied \textbf{before} broadening. Wavelength is
recomputed from the calibrated energy via
\(\lambda = 1239.841984\,\text{eV\,nm}\,/\,E_{\text{cal}}\) only when
the linear energy map is non-identity. The identity calibration
(\(a=1\), \(b=0\), \(\alpha=1\), no band overrides) produces
bit-identical output to running without \texttt{-{}-calibration}.

The user-selectable knob is:
\begin{itemize}
\tightlist
\item
  CLI: \texttt{-{}-calibration PATH} (path to a TOML file)
\item
  Library: \texttt{calibration=Calibration(\ldots)} parameter of
  \texttt{analyze()}
\end{itemize}

\textbf{Note.} The HTTP API does not accept calibration files. Use the
CLI or the Python library for this feature.

\begin{center}\rule{0.5\linewidth}{0.5pt}\end{center}

\hypertarget{toml-schema}{
\subsection{TOML schema}\label{toml-schema}}

\begin{Shaded}
\begin{Highlighting}[]
[energy]
a = 1.0    # multiplicative factor (default 1.0 = identity)
b = 0.0    # additive shift in eV (default 0.0)

[oscillator]
alpha = 1.0   # scaling factor for all oscillator strengths

[[band]]
index = 1           # 1-based excited-state index
sigma_ev = 0.25     # width override in eV (optional)
reorganization_ev = 0.40  # Marcus-Hush reorganization energy (optional)
\end{Highlighting}
\end{Shaded}

\begin{center}\rule{0.5\linewidth}{0.5pt}\end{center}

\hypertarget{what-you-will-produce}{
\subsection{What you will produce}\label{what-you-will-produce}}

\begin{longtable}[]{@{}lll@{}}
\toprule
Folder & By & Contents \\
\midrule
\endhead
\texttt{output/cli/tables/} & CLI & Result tables (CSV / JSON / XLSX) \\
\texttt{output/cli/plots/} & CLI & TIFF plots per sample, light source,
and (sample, light) pair \\
\bottomrule
\end{longtable}

Light sources used: \textbf{AM1.5G + LED-B4}.

\begin{center}\rule{0.5\linewidth}{0.5pt}\end{center}

\hypertarget{1-prerequisites}{
\subsection{1. Prerequisites}\label{1-prerequisites}}

If you have never run any case study before, follow
\href{https://github.com/pinarsyt/overlap-calculator/blob/main/case_studies/01_theoretical_only/README.md\#1-prerequisites}{Case
Study 01 \S1 Prerequisites} once to install the package and either the
Python or Docker option for the HTTP API.

\begin{center}\rule{0.5\linewidth}{0.5pt}\end{center}

\hypertarget{2-inputs}{
\subsection{2. Inputs}\label{2-inputs}}

The case-study \texttt{input.json} lists 2 Gaussian TD-DFT outputs. The
case study also includes a \texttt{calib.toml} file that leaves energy
and oscillator scaling at their identity values (\(a=1\), \(b=0\),
\(\alpha=1\)) and overrides the width of the first excited state to
\(\sigma = 0.25\)\,eV.

\hypertarget{3-run-with-the-cli}{
\subsection{3. Run with the CLI}\label{3-run-with-the-cli}}

\begin{Shaded}
\begin{Highlighting}[]
\ExtensionTok{overlap{-}calculator}\NormalTok{ analyze }\AttributeTok{{-}{-}input}\NormalTok{ case\_studies/14\_calibration\_block/input.json }\AttributeTok{{-}{-}out}\NormalTok{ case\_studies/14\_calibration\_block/output/cli }\AttributeTok{{-}{-}default{-}light{-}sources}\NormalTok{ AM15G,LEDB4 }\AttributeTok{{-}{-}calibration}\NormalTok{ case\_studies/14\_calibration\_block/calib.toml }\AttributeTok{{-}{-}plot{-}dpi}\NormalTok{ 400}
\end{Highlighting}
\end{Shaded}

The \texttt{run\_manifest.json} at the output root records the full
calibration block under \texttt{parameters.calibration}.

\begin{center}\rule{0.5\linewidth}{0.5pt}\end{center}

\hypertarget{4-http-api}{
\subsection{4. HTTP API}\label{4-http-api}}

The HTTP API does not accept calibration files. Use the CLI or the
Python library for this feature.

\begin{center}\rule{0.5\linewidth}{0.5pt}\end{center}

\hypertarget{5-drive-the-cli-from-one-python-script}{
\subsection{5. Drive the CLI from one Python
script}\label{5-drive-the-cli-from-one-python-script}}

\begin{Shaded}
\begin{Highlighting}[]
\ExtensionTok{python}\NormalTok{ case\_studies/14\_calibration\_block/submit.py}
\end{Highlighting}
\end{Shaded}

\begin{center}\rule{0.5\linewidth}{0.5pt}\end{center}

\hypertarget{6-expected-output-layout}{
\subsection{6. Expected output layout}\label{6-expected-output-layout}}

\begin{verbatim}
case_studies/14_calibration_block/output/
+-- cli/
    +-- run_manifest.json
    +-- tables/results.{csv,json,xlsx}        (2 samples x 2 light sources = 4 rows)
    +-- tables/results_timings.{csv,json,xlsx}
    +-- tables/descriptor_summary.{csv,xlsx}
    +-- tables/ranking_by_light_source__<metric>.{csv,json,xlsx}
    +-- tables/ranking_by_sample__<metric>.{csv,json,xlsx}
    +-- plots/
\end{verbatim}

Compare \texttt{gaussian\_lambda\_max\_nm} between the calibrated and
the default run: the \(-0.05\)\,eV shift should produce a red-shift in
the peak wavelength.

\begin{center}\rule{0.5\linewidth}{0.5pt}\end{center}

\hypertarget{7-related-case-studies}{
\subsection{7. Related case studies}\label{7-related-case-studies}}

\begin{itemize}
\tightlist
\item
  \href{https://github.com/pinarsyt/overlap-calculator/blob/main/case_studies/08_broadening_sigma/README.md}{Case
  Study 08 --- Broadening Sigma}
\item
  \href{https://github.com/pinarsyt/overlap-calculator/blob/main/case_studies/13_marcus_hush_width/README.md}{Case
  Study 13 --- Marcus--Hush Width}
\item
  \href{https://github.com/pinarsyt/overlap-calculator/blob/main/case_studies/15_vibronic_tabular_branch/README.md}{Case
  Study 15 --- Vibronic Tabular Branch}
\end{itemize}

\end{document}
